\documentclass[12pt,a4paper]{article}
\usepackage[utf8]{vietnam}
\usepackage{amsmath,amssymb}
\usepackage{tabularx}
\usepackage{tikz}
\usepackage{xcolor}
% Giảm lề tối đa tận dụng diện tích giấy (chuẩn máy in laser/photo)
\usepackage[top=0.6cm, bottom=0.6cm, left=0.6cm, right=0.6cm]{geometry}
\usepackage[dethi]{ex_test}

\setlength{\parskip}{0pt}

% Dòng kẻ chấm mờ đều đặn chuẩn viết tay (khoảng cách 7.2mm rộng rãi)
\newcommand{\dongke}[1]{%
  \begingroup%
  \setlength{\parskip}{0pt}%
  \baselineskip=20.5pt%
  \vspace{1.5mm}%
  \foreach \i in {1,...,#1}{%
    \noindent{\color{gray!60}\dotfill}\par%
  }%
  \vspace{1mm}%
  \endgroup%
}

% Định nghĩa ô bài tập tự luận trong bảng
\newcommand{\oBaiTap}[3][7]{%
  \vspace{1.5mm}%
  \noindent\textbf{#2} #3\par%
  \dongke{#1}%
}

% Câu trắc nghiệm 1 cột: 4 đáp án dàn hàng ngang trên 1 dòng
\newcommand{\cauTN}[6]{% \cauTN{số câu}{đề}{A}{B}{C}{D}
  \noindent\textbf{Câu #1.} #2\par\vspace{0.4mm}%
  \noindent%
  \makebox[0.25\linewidth][l]{\textbf{A.} #3}%
  \makebox[0.25\linewidth][l]{\textbf{B.} #4}%
  \makebox[0.25\linewidth][l]{\textbf{C.} #5}%
  \makebox[0.25\linewidth][l]{\textbf{D.} #6}\par\vspace{0.8mm}%
}

\begin{document}
\pagestyle{empty}

% #############################################################################
% ========================= TOÁN 6 - TUẦN 2 - BUỔI 2: MÃ ĐỀ 101 =============
% #############################################################################

% -------------------- TRANG 1 (MÃ ĐỀ 101): TRẮC NGHIỆM --------------------
\noindent
\begin{tabularx}{\linewidth}{|X|c|}
\hline
\textbf{BÀI 4: CỘNG VÀ TRỪ SỐ TỰ NHIÊN (TUẦN 2 - BUỔI 2)} & \textbf{Mã đề: 101} \\
\hline
\end{tabularx}

\vspace{1.2mm}

\noindent
\begin{tabularx}{\linewidth}{|X|}
\hline
\rule[-3mm]{0pt}{9.5mm}\textbf{Họ và tên:} \dotfill \quad \textbf{Lớp:} \makebox[1.8cm]{\dotfill} \\
\hline
\end{tabularx}

\vspace{1.2mm}

\noindent
\begin{tabularx}{\linewidth}{|p{3cm}|X|}
\hline
\centering\textbf{ĐIỂM SỐ} & \centering\textbf{LỜI NHẬN XÉT} \tabularnewline
\hline
\rule{0pt}{2.2cm} & \tabularnewline
\hline
\end{tabularx}

\vspace{1.5mm}
\noindent\textbf{BẢNG ĐIỀN ĐÁP ÁN TRẮC NGHIỆM (Mỗi câu đúng 0,25 điểm):}
\vspace{1.2mm}

\noindent
\renewcommand{\arraystretch}{1.5}
\begin{tabularx}{\linewidth}{|c|*{12}{>{\centering\arraybackslash}X|}}
\hline
\textbf{Câu} & \textbf{1} & \textbf{2} & \textbf{3} & \textbf{4} & \textbf{5} & \textbf{6} & \textbf{7} & \textbf{8} & \textbf{9} & \textbf{10} & \textbf{11} & \textbf{12} \\
\hline
\textbf{Đ/A} & & & & & & & & & & & & \\
\hline
\end{tabularx}
\renewcommand{\arraystretch}{1.0}

\vspace{3mm}

\noindent\textbf{\large PHẦN I. TRẮC NGHIỆM KHÁCH QUAN (3,0 điểm)}
\vspace{1.5mm}

\cauTN{1}{Trong phép tính $a + b = c$, các số $a$ và $b$ được gọi là:}{Số bị trừ}{Hiệu}{Số hạng}{Tổng}

\cauTN{2}{Điều kiện để thực hiện được phép trừ $a - b$ trong tập hợp số tự nhiên là:}{$a < b$}{$a \le b$}{$a \ge b$}{$a > b$}

\cauTN{3}{Tính chất $a + b = b + a$ thể hiện tính chất nào của phép cộng?}{Kết hợp}{Giao hoán}{Phân phối}{Cộng với số 0}

\cauTN{4}{Kết quả của phép tính nhẩm $138 + 262$ là:}{$390$}{$400$}{$410$}{$300$}

\cauTN{5}{Khẳng định nào sau đây là \textbf{sai}?}{$a + 0 = a$}{$a - 0 = a$}{$a - a = 0$}{$0 - a = a$}

\cauTN{6}{Quy tắc nào sau đây biểu diễn phép trừ một tổng: $a - (b + c) = $}{$a - b + c$}{$a - b - c$}{$a + b - c$}{$a + b + c$}

\cauTN{7}{Tìm số tự nhiên $x$, biết $x + 45 = 100$:}{$x = 45$}{$x = 55$}{$x = 65$}{$x = 145$}

\cauTN{8}{Tìm số tự nhiên $x$, biết $120 - x = 35$:}{$x = 85$}{$x = 95$}{$x = 155$}{$x = 75$}

\cauTN{9}{Tổng của $5$ số tự nhiên liên tiếp bắt đầu từ $1$ là:}{$10$}{$15$}{$20$}{$12$}

\cauTN{10}{Biết $(x - 12) + 20 = 50$, giá trị của $x$ là:}{$x = 42$}{$x = 38$}{$x = 18$}{$x = 62$}

\cauTN{11}{Buổi sáng bán $120\text{ kg}$ gạo, buổi chiều bán nhiều hơn sáng $30\text{ kg}$. Cả ngày bán được:}{$150\text{ kg}$}{$270\text{ kg}$}{$240\text{ kg}$}{$300\text{ kg}$}

\cauTN{12}{Tính nhanh tổng $19 + 28 + 81 + 72$ ta được kết quả là:}{$200$}{$190$}{$210$}{$180$}

\newpage
% -------------------- TRANG 2 (MÃ ĐỀ 101): TỰ LUẬN --------------------
\noindent
\begin{tabularx}{\linewidth}{|X|c|}
\hline
\textbf{PHẦN II. TỰ LUẬN (7,0 ĐIỂM)} & \textbf{Mã đề: 101} \\
\hline
\end{tabularx}

\vspace{1.2mm}

\noindent
\begin{tabularx}{\linewidth}{|X|}
\hline
\rule[-3mm]{0pt}{9.5mm}\textbf{Họ và tên:} \dotfill \quad \textbf{Lớp:} \makebox[1.8cm]{\dotfill} \\
\hline
\end{tabularx}

\vspace{2.5mm}

\noindent\textbf{\large Câu 1 (2,5 điểm). Tính bằng cách hợp lý (nếu có thể):}
\vspace{1.2mm}

\noindent
\begin{tabularx}{\linewidth}{|X|X|}
\hline
\oBaiTap[7]{a)}{$135 + 360 + 65 + 40$} & 
\oBaiTap[7]{b)}{$472 - (172 + 89)$} \\
\hline
\end{tabularx}

\vspace{3.5mm}

\noindent\textbf{\large Câu 2 (2,5 điểm). Tìm số tự nhiên $x$, biết:}
\vspace{1.2mm}

\noindent
\begin{tabularx}{\linewidth}{|X|X|}
\hline
\oBaiTap[7]{a)}{$x + 235 = 450 - 65$} & 
\oBaiTap[7]{b)}{$185 - (x - 15) = 90$} \\
\hline
\end{tabularx}

\vspace{3.5mm}

\noindent\textbf{\large Câu 3 (2,0 điểm). Bài toán thực tế:}
\vspace{1.2mm}

\noindent
\begin{tabularx}{\linewidth}{|X|}
\hline
\vspace{0.8mm}
Mẹ Lan đi siêu thị mua: $1$ bao gạo giá $185\,000$ đồng, $2$ vỉ trứng gà giá $65\,000$ đồng và $1$ chai dầu ăn giá $50\,000$ đồng. Mẹ Lan đưa cho nhân viên thu ngân một tờ tiền mệnh giá $500\,000$ đồng. Hỏi nhân viên thu ngân phải trả lại mẹ Lan bao nhiêu tiền thừa?
\par\vspace{1.5mm}
\noindent\textit{Bài giải:}
\dongke{8}
\\ \hline
\end{tabularx}

\newpage
% #############################################################################
% ========================= TOÁN 6 - TUẦN 2 - BUỔI 2: MÃ ĐỀ 102 =============
% #############################################################################

% -------------------- TRANG 3 (MÃ ĐỀ 102): TRẮC NGHIỆM --------------------
\noindent
\begin{tabularx}{\linewidth}{|X|c|}
\hline
\textbf{BÀI 4: CỘNG VÀ TRỪ SỐ TỰ NHIÊN (TUẦN 2 - BUỔI 2)} & \textbf{Mã đề: 102} \\
\hline
\end{tabularx}

\vspace{1.2mm}

\noindent
\begin{tabularx}{\linewidth}{|X|}
\hline
\rule[-3mm]{0pt}{9.5mm}\textbf{Họ và tên:} \dotfill \quad \textbf{Lớp:} \makebox[1.8cm]{\dotfill} \\
\hline
\end{tabularx}

\vspace{1.2mm}

\noindent
\begin{tabularx}{\linewidth}{|p{3cm}|X|}
\hline
\centering\textbf{ĐIỂM SỐ} & \centering\textbf{LỜI NHẬN XÉT} \tabularnewline
\hline
\rule{0pt}{2.2cm} & \tabularnewline
\hline
\end{tabularx}

\vspace{1.5mm}
\noindent\textbf{BẢNG ĐIỀN ĐÁP ÁN TRẮC NGHIỆM (Mỗi câu đúng 0,25 điểm):}
\vspace{1.2mm}

\noindent
\renewcommand{\arraystretch}{1.5}
\begin{tabularx}{\linewidth}{|c|*{12}{>{\centering\arraybackslash}X|}}
\hline
\textbf{Câu} & \textbf{1} & \textbf{2} & \textbf{3} & \textbf{4} & \textbf{5} & \textbf{6} & \textbf{7} & \textbf{8} & \textbf{9} & \textbf{10} & \textbf{11} & \textbf{12} \\
\hline
\textbf{Đ/A} & & & & & & & & & & & & \\
\hline
\end{tabularx}
\renewcommand{\arraystretch}{1.0}

\vspace{3mm}

\noindent\textbf{\large PHẦN I. TRẮC NGHIỆM KHÁCH QUAN (3,0 điểm)}
\vspace{1.5mm}

\cauTN{1}{Trong phép tính $m - n = p$, số $n$ được gọi là:}{Số bị trừ}{Số trừ}{Hiệu}{Số hạng}

\cauTN{2}{Phép trừ nào sau đây \textbf{không} thực hiện được trong tập hợp số tự nhiên $\mathbb{N}$?}{$15 - 8$}{$20 - 20$}{$7 - 12$}{$100 - 99$}

\cauTN{3}{Tính chất $(a + b) + c = a + (b + c)$ thể hiện tính chất nào của phép cộng?}{Giao hoán}{Kết hợp}{Cộng với số 0}{Phân phối}

\cauTN{4}{Kết quả của phép tính nhẩm $245 + 155$ là:}{$400$}{$390$}{$410$}{$300$}

\cauTN{5}{Khẳng định nào sau đây là \textbf{đúng}?}{$x - x = x$}{$x + 0 = 0$}{$x - 0 = x$}{$0 - x = x$}

\cauTN{6}{Quy tắc nào sau đây biểu diễn phép trừ một tổng: $m - (n + p) = $}{$m - n + p$}{$m + n - p$}{$m - n - p$}{$m + n + p$}

\cauTN{7}{Tìm số tự nhiên $y$, biết $y + 35 = 90$:}{$y = 45$}{$y = 55$}{$y = 65$}{$y = 125$}

\cauTN{8}{Tìm số tự nhiên $y$, biết $150 - y = 45$:}{$y = 95$}{$y = 105$}{$y = 115$}{$y = 195$}

\cauTN{9}{Tổng của $5$ số tự nhiên liên tiếp bắt đầu từ $0$ là:}{$10$}{$15$}{$5$}{$20$}

\cauTN{10}{Biết $(x - 15) + 30 = 70$, giá trị của $x$ là:}{$x = 55$}{$x = 25$}{$x = 85$}{$x = 45$}

\cauTN{11}{Buổi sáng mượn $110$ cuốn sách, buổi chiều mượn nhiều hơn sáng $25$ cuốn. Cả ngày mượn được:}{$135\text{ cuốn}$}{$245\text{ cuốn}$}{$220\text{ cuốn}$}{$235\text{ cuốn}$}

\cauTN{12}{Tính nhanh tổng $26 + 37 + 74 + 63$ ta được kết quả là:}{$190$}{$200$}{$210$}{$180$}

\newpage
% -------------------- TRANG 4 (MÃ ĐỀ 102): TỰ LUẬN --------------------
\noindent
\begin{tabularx}{\linewidth}{|X|c|}
\hline
\textbf{PHẦN II. TỰ LUẬN (7,0 ĐIỂM)} & \textbf{Mã đề: 102} \\
\hline
\end{tabularx}

\vspace{1.2mm}

\noindent
\begin{tabularx}{\linewidth}{|X|}
\hline
\rule[-3mm]{0pt}{9.5mm}\textbf{Họ và tên:} \dotfill \quad \textbf{Lớp:} \makebox[1.8cm]{\dotfill} \\
\hline
\end{tabularx}

\vspace{2.5mm}

\noindent\textbf{\large Câu 1 (2,5 điểm). Tính bằng cách hợp lý (nếu có thể):}
\vspace{1.2mm}

\noindent
\begin{tabularx}{\linewidth}{|X|X|}
\hline
\oBaiTap[7]{a)}{$142 + 250 + 58 + 50$} & 
\oBaiTap[7]{b)}{$563 - (163 + 78)$} \\
\hline
\end{tabularx}

\vspace{3.5mm}

\noindent\textbf{\large Câu 2 (2,5 điểm). Tìm số tự nhiên $x$, biết:}
\vspace{1.2mm}

\noindent
\begin{tabularx}{\linewidth}{|X|X|}
\hline
\oBaiTap[7]{a)}{$x + 145 = 380 - 75$} & 
\oBaiTap[7]{b)}{$195 - (x - 25) = 80$} \\
\hline
\end{tabularx}

\vspace{3.5mm}

\noindent\textbf{\large Câu 3 (2,0 điểm). Bài toán thực tế:}
\vspace{1.2mm}

\noindent
\begin{tabularx}{\linewidth}{|X|}
\hline
\vspace{0.8mm}
Bác An đi nhà sách mua: $1$ bộ sách giáo khoa giá $195\,000$ đồng, $2$ tập vở ô ly giá $75\,000$ đồng và $1$ hộp bút sáp màu giá $40\,000$ đồng. Bác An đưa cho nhân viên thu ngân một tờ tiền mệnh giá $500\,000$ đồng. Hỏi nhân viên thu ngân phải trả lại bác An bao nhiêu tiền thừa?
\par\vspace{1.5mm}
\noindent\textit{Bài giải:}
\dongke{8}
\\ \hline
\end{tabularx}

\end{document}
